\section{Conclusion}
\label{sec:conclusion}
This paper presented RAVM-DT for continual maintenance of deployed DT models under evolving and potentially unreliable monitoring data. By separating candidate adaptation from deployment commitment, RAVM-DT uses paired recent and bounded historical verification evidence to decide whether an adapted candidate should replace the current state. Across the five primary lifecycle conditions, it reduces harmful-transition frequency by about 85.5\%--91.4\% and harmful downside by 85.3\%--94.9\% relative to active adaptation baselines, at a 6.8\%--13.2\% MSE increase; relative to Frozen, it lowers MSE by about 34.9\% while retaining nonzero adaptation. Mechanism, benign-drift, and architecture/horizon results support the intended selective-transition behavior rather than a retain-all policy. On Alibaba workloads, the same qualitative safety control transfers, but the 6.49\% mixed-condition harmful-transition frequency exceeds the 5\% synthetic calibration criterion, indicating domain-dependent risk calibration. Current evidence is empirical and limited to workload prediction on one controlled benchmark and one external trace; future work will study calibration under domain shift, additional DT tasks/platforms, and mechanisms that improve adaptation without sacrificing transition safety.
